% Удиртгал (дугааргүй)
% Шинэчилсэн: 2026-10-05. Эх template: ~/Downloads/template/Chapters/Introduction.tex
% БИЧИХ ДАРААЛАЛ (багшийн заавар): 1) Судлах үндэслэл -> 2) Судлагдсан байдал ->
%   ... -> хамгийн сүүлд Зорилго, Зорилт.  Гарчгийн дараалал баримтад доорх шиг байна.

\phantomsection
\addchaptertocentry{Удиртгал}
\label{Introduction}

\section*{Удиртгал}
% Ноорог v1 — 2026-10-05
Програм хангамжийн шаардлага болон хичээлийн материал нь хоёулаа байгалийн хэлээр бичигддэг баримт бөгөөд тэдгээрийн чанар нь дараагийн үйл ажиллагааны үр дүнг тодорхойлдог. Шаардлага бүрхэг эсвэл давхардсан бол хөгжүүлэгч, захиалагч хоёр түүнийг өөр өөрөөр ойлгож, алдаа хожуу илэрдэг. Үүнтэй адил хичээлийн суралцахуйн зорилт лекц, даалгавартаа тусгагдаагүй бол суралцагч тухайн зорилтод хүрэх боломж хязгаарлагдана. Ийм баримтыг хянах ажил одоогоор ихэвчлэн гараар хийгддэг бөгөөд монгол хэл дээр түүнийг автоматжуулах хэрэгсэл ховор байна. Энэхүү ажлаар хиймэл оюун ухаан болон байгалийн хэлний боловсруулалтын аргуудыг ашиглан монгол хэл дээрх эдгээр баримтын чанарыг үнэлэх систем хөгжүүлж, аргуудыг харьцуулан судална.

\subsection*{Судлах үндэслэл}
% Ноорог v1 — 2026-10-05. [..] тэмдэглэгээтэй хэсгүүдийг өгөгдөл цуглуулсны дараа нөхнө.
Програм хангамжийн шаардлага (software requirement) нь захиалагч, хэрэглэгч болон хөгжүүлэгчийн хооронд тохиролцсон системийн үйл ажиллагаа, шинж чанарыг тодорхойлдог тул түүний чанар нь зохиомж, хөгжүүлэлт, туршилт зэрэг дараагийн бүх шатны үр дүнд шууд нөлөөлдөг. Шаардлагын инженерчлэлийн олон улсын ISO/IEC/IEEE 29148 стандарт нь шаардлага бүр нэг утгатай (unambiguous), шалгагдахуйц (verifiable), бүрэн (complete) байх ёстойг тодорхойлсон \cite{iso29148}. Өөрөөр хэлбэл, \enquote{систем хурдан ажиллана}, \enquote{хэрэглэхэд хялбар байна} гэх мэт хэмжих шалгуургүй шаардлагыг биелсэн эсэхийг бодитойгоор шалгах боломжгүй юм.

Шаардлагын алдааг хожуу илрүүлэх нь их зардал дагуулдаг. Боэм, Басили нарын тоймд програм хангамжийн алдааг хэрэглэгчид хүлээлгэн өгсний дараа илрүүлж засах нь шаардлага, зохиомжийн шатанд засахаас ихэвчлэн 100 дахин их зардалтай гэж дүгнэсэн бөгөөд жижиг, эгзэгтэй бус системийн хувьд энэ харьцаа 5:1 орчим байдгийг мөн тэмдэглэсэн \cite{boehm2001}. Иймд шаардлагын чанарыг баримт бичгийн шатанд нь шалгах нь зардлыг бууруулах үр дүнтэй арга юм.

Монгол Улсад төрийн худалдан авах ажиллагааны цахим систем (tender.gov.mn)-ээр зарлагдаж буй тендерийн баримт бичигт ажлын даалгавар, техникийн тодорхойлолт хавсаргагддаг. Төрийн болон орон нутгийн өмчийн хөрөнгөөр бараа, ажил, үйлчилгээ худалдан авах тухай хуулийн 12.2-т техникийн тодорхойлолтод шалгуур үзүүлэлтийг тодорхойлохыг заасан байдаг \cite{huuli2023}. Гэвч програм хангамжийн худалдан авалтын ажлын даалгаварт \enquote{хурдан}, \enquote{хялбар}, \enquote{шаардлагатай бол} гэх мэт хэмжих шалгуургүй, бүрхэг илэрхийлэл болон давхардсан шаардлага ажиглагдаж байна [бодит жишээ 2--3: баримтын дугаар, огноо; урьдчилсан тоо — өгөгдөл цуглуулсны дараа нөхнө]. Ийм баримтыг гараар хянах нь цаг их шаарддаг бөгөөд хянагчийн туршлагаас хамаарч үр дүн нь өөр өөр гарч болно.

Шаардлагын баримтыг байгалийн хэлний боловсруулалт (Natural Language Processing, NLP)-ын аргаар автоматаар шинжлэх чиглэл буюу NLP4RE (NLP for Requirements Engineering) нь олон улсад өргөн судлагдаж байна. Жао нарын системчилсэн зураглалын судалгаагаар энэ чиглэлийн 404 судалгааг тодорхойлсон ч тэдгээрийн зөвхөн 7\% нь бодит орчинд үнэлэгдсэн, илрүүлсэн 130 хэрэгслийн 13.08\% нь л татаж авах боломжтой байсан \cite{zhao2021}. Харин монгол хэл дээрх шаардлагын баримтыг автоматаар шалгах хэрэгсэл, судалгаа бидний хайлтаар олдоогүй [хайсан сан, түлхүүр үг, хугацааг заана]. Монгол хэл нь залгамал бүтэцтэй, үгийн хувилгавар ихтэй, мөн тэмдэглэгээтэй өгөгдөл цөөн бага нөөцтэй (low-resource) хэл тул англи хэлэнд зориулсан аргуудыг шууд хэрэглэх боломж хязгаарлагдмал [эх сурвалж хэрэгтэй].

Үүнтэй төстэй асуудал боловсролын салбарт ч тулгардаг. Хичээлийн суралцахуйн зорилт, сургалтын үйл ажиллагаа, үнэлгээ хоорондоо уялдсан байх нь сургалтын чанарын үндсэн зарчмын нэг гэж үздэг \cite{biggs1996}. Хичээлийн хөтөлбөрт тусгагдсан зорилт лекц, даалгаварт хэрхэн хамрагдсаныг гараар шалгах нь мөн л баримт бичгийг нарийвчлан харьцуулах ажил бөгөөд утгын ижил төстэй байдлыг тооцох ижил аргуудаар автоматжуулах боломжтой.

Дээрх үндэслэлээр хиймэл оюун ухааны аргуудыг ашиглан монгол хэл дээрх програм хангамжийн шаардлагын бүрхэг, давхардсан байдлыг илрүүлэх, мөн хичээлийн материалын уялдааг шалгах системийг хөгжүүлж, бага нөөцтэй нөхцөлд аль арга илүү тохиромжтойг харьцуулан судлах шаардлагатай гэж үзлээ.

\subsection*{Судлагдсан байдал}
% Ноорог v1 — 2026-10-05. Энд товч тойм; дэлгэрэнгүй шинжилгээ, харьцуулсан хүснэгтийг 1.6-р хэсэгт.
\paragraph{Олон улсад.}
Шаардлагын баримтын чанарыг автоматаар шалгах судалгаа эрт эхэлсэн бөгөөд эхний үеийн аргууд нь бүрхэг үг, хэллэгийн жагсаалт болон хэлзүйн дүрэмд суурилж байв. Энэ чиглэлийн төлөөлөх ажлын нэг болох Фэммер нарын судалгаа нь программ кодын \enquote{муу үнэр} (code smell)-ийн санааг шаардлагад шилжүүлж, субъектив хэллэг, бүрхэг дайвар үг ба тэмдэг нэр, хэмжигдэхгүй нээлттэй нэр томьёо, давуу ба харьцуулсан зэрэг, үгүйсгэл, бүрхэг төлөөний үг зэрэг найман төрлийн \enquote{шаардлагын үнэр} (requirements smell)-ийг тодорхойлж, Smella хэрэгслийг хөгжүүлсэн. Аж үйлдвэрийн болон академийн баримтууд дээр уг хэрэгсэл дунджаар 59\%-ийн нарийвчлал (precision), 82\%-ийн бүрэн илрүүлэлт (recall)-тэй ажилласан бөгөөд үнэрийн төрлөөс хамаарч нарийвчлал ихээхэн ялгаатай байсан \cite{femmer2017}. Энэ нь дүрэмд суурилсан арга хурдан, тайлбарлагдахуйц боловч хуурамч илрүүлэлт ихтэй байж болохыг харуулж байна.

Давхардсан буюу утгаараа ижил шаардлагыг илрүүлэх асуудлыг Фалесси нар судалж, байгалийн хэлний боловсруулалтын олон аргыг хослуулан аж үйлдвэрийн нөхцөлд үнэлсэн \cite{falessi2013}. Сүүлийн жилүүдэд текстийн embedding болон Transformer загварын хөгжлөөр утгын ижил төстэй байдлыг тооцох аргууд түгээмэл болж, улмаар их хэлний загвар (Large Language Model, LLM)-ыг шаардлагын инженерчлэлд ашиглах судалгаа эрчимжиж байна. Тухайлбал, Башир нарын аж үйлдвэрийн судалгаагаар харьцангуй жижиг, нээлттэй LLM-үүдэд арван жишээ (10-shot) өгөхөд бүрхэг шаардлагыг ангилах гүйцэтгэл жишээгүй (0-shot) үеийнхээс дунджаар 20.2\%-иар нэмэгдсэн ч салбарын тусгай нэр томьёог ойлгох чадвар хязгаарлагдмал байсныг тэмдэглэсэн \cite{bashir2025}. Мөн судлаачдад нийтэд нээлттэй шаардлагын баримт хүртээмжтэй байлгах зорилгоор 79 баримт бүхий PURE өгөгдлийн багц бүрдүүлэгдсэн \cite{ferrari2017pure}. Гэвч Жао нарын тоймд дурдсанчлан энэ чиглэлийн ихэнх судалгаа лабораторийн туршилтаар хязгаарлагдаж, хэрэгслүүдийн цөөн хувь нь л нээлттэй байна \cite{zhao2021}. Дээрх судалгаанууд нь үндсэндээ англи хэл дээрх баримтад хийгдсэн [нотлох эсэхийг 1.6-р хэсэгт шалгана].

Хичээлийн материалын уялдааны талаар суралцахуйн зорилт, сургалтын үйл ажиллагаа, үнэлгээг уялдуулах онолын суурийг Биггс тавьсан бол \cite{biggs1996}, суралцахуйн зорилтыг танин мэдэхүйн түвшнээр ангилах нийтлэг хүрээ нь Блумын шинэчилсэн ангилал (revised Bloom's taxonomy) юм \cite{anderson2001}. Ли нарын судалгаанд 5,558 хичээлийн 21,380 суралцахуйн зорилтыг Блумын түвшнээр автоматаар ангилахад BERT загвар F1 оноо 0.95, Cohen's kappa 0.93 хүртэл хүрсэн бөгөөд SVM, санамсаргүй ой (random forest), XGBoost зэрэг уламжлалт машин сургалтын аргууд түүнтэй ойролцоо үр дүн үзүүлсэн \cite{li2022bloom}. Заки нар хичээлийн суралцахуйн үр дүнг хөтөлбөрийн үр дүнтэй BERT болон байгалийн хэлний боловсруулалтын аргаар автоматаар холбож, хоёр хөтөлбөрийн бодит өгөгдөл дээр мэргэжилтний үнэлгээтэй харьцуулахад 83.1\% ба 88.1\%-ийн нарийвчлалтай ажилласан бөгөөд вэб хэрэглээ хөгжүүлж, өгөгдөл, кодоо нээлттэй болгосон \cite{zaki2023}. Эдгээр судалгаа нь зорилтыг ангилах, эсвэл хичээлийн болон хөтөлбөрийн түвшний зорилтыг хооронд нь холбоход төвлөрсөн бөгөөд нэг хичээлийн дотор зорилт, лекц, даалгаврын гурвалсан уялдааг шалгах, хамрагдаагүй зорилт болон зорилгогүй даалгаврыг илрүүлэх асуудал [1.6-р хэсэгт нэмэлт эх сурвалжаар нягтлах] энэхүү ажлын Хэсэг Б-ийн гол чиглэл болно.

\paragraph{Монголд.}
Монгол хэл нь олон хэлт загваруудын сургалтын өгөгдөлд маш бага хэмжээгээр төлөөлөгддөг. Жишээлбэл, XLM-RoBERTa (XLM-R) загварыг сургахад ашигласан CC-100 корпуст монгол хэлний өгөгдөл 248 сая токен (3.0 GiB) байхад англи хэлнийх 55,608 сая токен (300.8 GiB) буюу ойролцоогоор 224 дахин их байна \cite{conneau2020}. Кирилл үсгийн монгол хэлэнд зориулсан BERT загварыг Википедиа болон мэдээний корпус дээр сургаж нээлттэй байршуулсан ажил байдаг \cite{mongolianbert}. [Монгол хэлний үгзүйн шинжилгээ, текст ангиллын судалгаа 1--2, эх сурвалжтай.] Монгол хэлний боловсруулалтын судалгааг шүүхдээ уламжлалт монгол бичиг болон кирилл бичгийн судалгааг ялгах шаардлагатай.

Харин монгол хэл дээрх програм хангамжийн шаардлага, тендерийн ажлын даалгаврын чанарыг автоматаар үнэлсэн судалгаа, хэрэгсэл бидний хайлтаар олдоогүй [хайсан сан: Google Scholar, ResearchGate, ШУТИС-ийн номын сан; түлхүүр үг; хайсан огноо]. ШУТИС болон бусад их сургуулийн ижил төстэй дипломын ажлууд [нэр, он, сэдэв, ялгаа — 2--3 ажил] нь шаардлагын шинжилгээг гараар хийсэн бөгөөд тэдгээрийн \enquote{шаардлагын шинжилгээ} бүлгүүдийг зөвшөөрөлтэйгээр энэхүү ажлын үнэлгээний өгөгдөлд ашиглана.

\subsection*{Судалгааны асуулт}
% Ноорог v1 — 2026-10-05. СА-Б-г О.Мэнд-Амар нягтална.
Энэхүү ажлын хүрээнд дараах судалгааны асуултад хариулна.

\begin{enumerate}[nosep]
    \item \textbf{СА-А1.} Бага нөөцтэй монгол хэл дээр (1) дүрэмд суурилсан, (2) embedding болон fine-tune хийсэн жижиг загвар (XLM-R), (3) их хэлний загвар (LLM), (4) эдгээрийг хосолсон аргын аль нь бүрхэг (хэмжих боломжгүй) болон давхардсан шаардлагыг хамгийн сайн илрүүлэх вэ?
    \item \textbf{СА-А2.} Дээрх аргууд нарийвчлал, тогтвортой байдал, боловсруулах хурд, зардлын хувьд хэрхэн ялгаатай вэ?
    \item \textbf{СА-Б.} Embedding-д суурилсан утгын ижил төстэй байдал болон LLM-д суурилсан арга нь монгол хэл дээрх хичээлийн суралцахуйн зорилт, лекц, даалгаврын хоорондын холбоог мэргэжилтний тэмдэглэгээтэй хэр нийцтэй тодорхойлох вэ?
\end{enumerate}

Аргуудыг үнэлэхдээ ангиллын тэнцвэргүй байдлыг харгалзан accuracy-г бус, нарийвчлал (Precision), бүрэн илрүүлэлт (Recall), F1 оноог ашиглана. Тэмдэглэгээний найдвартай байдлыг Cohen's kappa коэффициентоор, LLM-ийн тогтвортой байдлыг ижил оролтоор давтан ажиллуулж гарсан үр дүнгийн зөрүүгээр хэмжинэ. Мөн боловсруулах хугацаа болон API-ийн зардлыг тооцно.

\subsection*{Шинэлэг тал}
% Ноорог v1 — 2026-10-05. "бидний мэдэх хүрээнд" гэсэн болгоомжлолыг хадгална.
Бидний мэдэх хүрээнд энэхүү ажлын шинэлэг тал нь дараах байдлаар илэрнэ.
\begin{enumerate}[nosep]
    \item ISO/IEC/IEEE 29148 стандартад тулгуурласан тэмдэглэгээний дүрмээр бэлтгэсэн, монгол хэл дээрх програм хангамжийн шаардлагын чанарын тэмдэглэгээтэй өгөгдлийн багцыг бүрдүүлж, тэмдэглэгээний тохироог хэмжинэ.
    \item Дүрэм, жижиг загвар, LLM болон хосолсон аргыг монгол хэл дээрх ижил өгөгдөл, ижил нөхцөлд харьцуулж, бага нөөцтэй хэлэнд аль арга илүү тохиромжтойг тоон үзүүлэлтээр харуулна.
    \item LLM-ийн үр дүнгийн тогтвортой байдлыг давтан туршилтаар хэмжиж, бусад аргатай харьцуулна.
    \item Нэгдсэн шалгагч (Checker) архитектураар програм хангамжийн шаардлага болон хичээлийн материал гэсэн хоёр төрлийн баримтыг нэг системд шалгана.
\end{enumerate}

\subsection*{Практик ач холбогдол}
% Ноорог v1 — 2026-10-05
Хөгжүүлсэн систем нь дараах хэрэглэгчдэд ашиг тустай.
\begin{itemize}[nosep]
    \item \textbf{Тендерийн баримт бэлтгэгч, хянагч:} ажлын даалгаврыг зарлахаас өмнө бүрхэг, давхардсан шаардлагыг илрүүлж, засах санал авах;
    \item \textbf{Програм хангамжийн баг, бизнес шинжээч:} шаардлагын баримтыг хөгжүүлэлт эхлэхээс өмнө хянах;
    \item \textbf{Оюутан:} бие даалт, дипломын ажлын шаардлагын шинжилгээ бүлгийн чанарыг өөрөө шалгах;
    \item \textbf{Багш, хөтөлбөрийн баг:} хичээлийн зорилт лекц, даалгаварт хэрхэн хамрагдсаныг матрицаар харж, дутуу хамрагдсан зорилт болон зорилгогүй даалгаврыг илрүүлэх.
\end{itemize}
Систем нь шийдвэрийг хүний оронд гаргахгүй, харин хянагчид анхаарах ёстой мөрийг тайлбартай нь санал болгох туслах хэрэгсэл юм.

\subsection*{Хамрах хүрээ ба хязгаарлалт}
% Ноорог v1 — 2026-10-05
\textbf{Ажлын хүрээнд хийх зүйлс:}
\begin{itemize}[nosep]
    \item \textbf{Хэсэг А:} монгол хэл дээрх шаардлагаас бүрхэг (хэмжих боломжгүй) болон давхардсан шаардлагыг илрүүлэх, илрүүлэлт бүрт тайлбар, засах санал гаргах. Цаг хугацаа үлдвэл хоорондоо зөрчилдсөн шаардлагыг илрүүлэх (NLI, LLM).
    \item \textbf{Хэсэг Б:} хичээлийн суралцахуйн зорилт, лекц, даалгаврын холбоог тодорхойлох, хамрагдаагүй зорилт болон зорилгогүй даалгаврыг илрүүлэх, хамрах хүрээний матриц (heatmap) гаргах.
    \item Дээрх шалгалтуудыг вэб системээр (React, FastAPI, PostgreSQL) хэрэглэгчид хүргэх.
\end{itemize}

\textbf{Ажлын хүрээнд хийхгүй зүйлс:} шаардлагыг автоматаар бичих; шаардлагын мөшгөлт (traceability)-ийн бүрэн систем; скандсан PDF баримтаас текст таних (OCR); баримтыг бүрэн автоматаар шаардлага болгон хуваах; төслийн удирдлагын систем.

\textbf{Хязгаарлалт:} Үнэлгээний өгөгдөл нь tender.gov.mn-ийн нээлттэй ажлын даалгавар болон зөвшөөрөлтэйгөөр авсан өмнөх дипломын ажлуудаас бүрдсэн 150--200 шаардлагаар хязгаарлагдана. Сургалтын өгөгдөл нь англи хэлнээс орчуулж гараар шалгасан болон бодит шаардлагаас үүсгэсэн 500--1,000 шаардлага байх тул орчуулгын хэв маяг загварт нөлөөлөх эрсдэлтэй. Шаардлагын бүрхэг эсэх нь тодорхой хэмжээгээр субъектив тул тэмдэглэгээний дүрмийг давтан сайжруулж, тохироог kappa-гаар хэмжинэ. Систем нь текст хэлбэрт хөрвүүлэх боломжтой баримттай ажиллана.

\subsection*{Ажлын хуваарилалт}
% Ноорог v1 — 2026-10-05. Declaration-ийн шаардлагаар хэн юу хийснийг ил тод бичнэ.
Энэхүү ажлыг хоёр гүйцэтгэгч хамтран хийсэн бөгөөд хуваарилалтыг Хүснэгт \ref{tab:huvaarilalt}-д харуулав.

\begin{table}[H]
\centering
\caption{Гүйцэтгэгчдийн ажлын хуваарилалт}
\label{tab:huvaarilalt}
\begin{tabular}{|p{3.2cm}|p{5.6cm}|p{5.6cm}|}
\hline
\textbf{Хэсэг} & \textbf{Г.Нармандах} & \textbf{О.Мэнд-Амар} \\ \hline
Судалгааны чиглэл & Хэсэг А: шаардлагын чанар (бүрхэг, давхардсан, нэмэлтээр зөрчилдсөн шаардлага) & Хэсэг Б: хичээлийн материалын уялдаа (зорилт, лекц, даалгавар) \\ \hline
Хөгжүүлэлт & Backend (FastAPI), шалгагчид, загвар сургалт & Frontend (React), хамрах хүрээний матриц, Docker байршуулалт \\ \hline
Өгөгдөл & Шаардлагын өгөгдөл цуглуулах, тэмдэглэх & Хичээлийн материалын өгөгдөл цуглуулах, тэмдэглэх \\ \hline
Хамтарсан & \multicolumn{2}{p{11.2cm}|}{Системийн архитектур, Checker интерфэйс, embedding үйлчилгээ, LLM адаптер, өгөгдлийн сан, удиртгал, онолын болон архитектурын бүлгүүд} \\ \hline
\end{tabular}
\end{table}

\subsection*{Зорилго}
% Ноорог v1 — 2026-10-05
Энэхүү ажлын зорилго нь хиймэл оюун ухаанд суурилсан аргуудаар монгол хэл дээрх програм хангамжийн шаардлагын бүрхэг болон давхардсан байдлыг илрүүлж, хичээлийн материалын уялдааг шалгах систем хөгжүүлэх, мөн бага нөөцтэй нөхцөлд аль арга илүү тохиромжтойг харьцуулан тогтооход оршино.

\subsection*{Зорилт}
Уг зорилгод хүрэхийн тулд дараах зорилтуудыг дэвшүүлж байна. Үүнд:
\begin{enumerate}[nosep]
    \item Шаардлагын чанар, хичээлийн уялдаа, байгалийн хэлний боловсруулалтын аргууд болон ижил төстэй судалгааг судлах (Бүлэг \ref{Chapter1});
    \item Монгол хэл дээрх шаардлага болон хичээлийн материалын өгөгдөл цуглуулж, ISO/IEC/IEEE 29148-д тулгуурласан тэмдэглэгээний дүрмээр тэмдэглэх;
    \item Системийн шаардлагыг тодорхойлж, шинжилгээний загваруудыг гаргах (Бүлэг 2);
    \item Давхаргат архитектур болон нийтлэг Checker интерфэйс бүхий системийг зохиомжилж, хөгжүүлэх (Бүлэг 3);
    \item Дүрэм, жижиг загвар, LLM болон хосолсон аргыг хэрэгжүүлж, ижил нөхцөлд туршин харьцуулах (Бүлэг 4);
    \item Туршилтын үр дүнд дүгнэлт хийж, аргуудын давуу ба сул талыг тодорхойлох.
\end{enumerate}

\subsection*{Дипломын бүтэц}
Дипломын ажил удиртгал, дөрвөн бүлэг, ерөнхий дүгнэлт болон хавсралтаас бүрдэнэ. Бүлэг 1-д шаардлагын чанар, хичээлийн материалын уялдааны онол, байгалийн хэлний боловсруулалтын аргууд, ижил төстэй судалгаа болон ашиглах технологийг судална. Бүлэг 2-т системийн шинжилгээ, Бүлэг 3-т системийн зохиомж ба хөгжүүлэлтийг тайлбарлана. Бүлэг 4-т өгөгдөл, тэмдэглэгээ, туршилтын арга зүй болон аргуудын харьцуулсан үр дүнг танилцуулна. Хавсралтад тэмдэглэгээний дүрэм, prompt-ууд болон чухал кодын хэсгүүдийг оруулна.
